% Chapter 4: Unsupervised Learning
\chapter{Unsupervised Learning}

\section{Objectives}
The unsupervised learning phase aims to discover hidden patterns not captured by supervised classification:
\begin{itemize}
    \item Identify geographic fire hotspots
    \item Group fires by environmental characteristics
    \item Define fire risk zones for management
\end{itemize}

\section{Feature Selection for Clustering}
Features selected based on relevance for pattern discovery:
\begin{itemize}
    \item \textbf{Geographic:} Latitude, longitude
    \item \textbf{Climate:} Drought index, summer heat, precipitation
    \item \textbf{Topographic:} Elevation, terrain complexity, coastal proximity
    \item \textbf{Soil:} Organic carbon, fine fraction
\end{itemize}

\textbf{Scaling decision:} StandardScaler applied to all features. Geographic coordinates scaled separately to preserve spatial relationships while ensuring equal contribution from all feature categories.

\section{K-Means Clustering}

\subsection*{Algorithm Choice}
K-Means selected because: (1) efficient for large datasets, (2) well-suited for convex clusters, (3) simple interpretation of centroids.

\subsection*{Determining Optimal K}
Two methods used to select the number of clusters:

\textbf{Elbow Method:} Computed within-cluster sum of squares (inertia) for $k = 2$ to 15. The "elbow" indicates diminishing returns from additional clusters.

\textbf{Silhouette Analysis:} Measures how similar points are to their own cluster versus others. Higher scores indicate better-defined clusters.

\begin{figure}[H]
    \centering
    \fbox{\parbox{0.85\textwidth}{\centering\vspace{3cm}\textbf{[PLACEHOLDER: Elbow Plot and Silhouette Scores]}\vspace{3cm}}}
    \caption{Cluster number selection: elbow method and silhouette analysis}
    \label{fig:kmeans_selection}
\end{figure}

\subsection*{Results}
K-Means clustering objective function (minimized):
\begin{equation}
    J = \sum_{k=1}^{K} \sum_{x_i \in C_k} ||x_i - \mu_k||^2
\end{equation}

\begin{figure}[H]
    \centering
    \fbox{\parbox{0.85\textwidth}{\centering\vspace{3cm}\textbf{[PLACEHOLDER: K-Means Spatial Cluster Map]}\vspace{3cm}}}
    \caption{Geographic distribution of K-Means clusters}
    \label{fig:kmeans_map}
\end{figure}

\textbf{Cluster interpretation:} Distinct fire zones corresponding to geographic and climate regions. Clear separation between coastal Mediterranean zones and interior areas.

\section{DBSCAN Clustering}

\subsection*{Algorithm Choice}
DBSCAN (Density-Based Spatial Clustering of Applications with Noise) selected because: (1) discovers arbitrary-shaped clusters, (2) automatically determines cluster count, (3) explicitly handles noise/outliers.

\subsection*{Parameter Selection}
\begin{itemize}
    \item \textbf{Epsilon ($\epsilon$):} Selected using k-distance plot---the knee point indicates appropriate neighborhood radius
    \item \textbf{Min\_samples:} Rule of thumb $\geq$ dimensionality + 1; tested range \{5, 10, 15, 20\}
\end{itemize}

\begin{figure}[H]
    \centering
    \fbox{\parbox{0.85\textwidth}{\centering\vspace{3cm}\textbf{[PLACEHOLDER: K-Distance Plot for Epsilon Selection]}\vspace{3cm}}}
    \caption{K-distance plot for DBSCAN parameter selection}
    \label{fig:dbscan_kdist}
\end{figure}

\subsection*{Results}
DBSCAN identifies:
\begin{itemize}
    \item \textbf{Core points:} Points with $\geq$ min\_samples neighbors within $\epsilon$
    \item \textbf{Border points:} Within $\epsilon$ of core but not core themselves
    \item \textbf{Noise points:} Isolated fire events (neither core nor border)
\end{itemize}

\begin{figure}[H]
    \centering
    \fbox{\parbox{0.85\textwidth}{\centering\vspace{3cm}\textbf{[PLACEHOLDER: DBSCAN Spatial Map with Noise Points]}\vspace{3cm}}}
    \caption{DBSCAN clusters with noise points shown in gray}
    \label{fig:dbscan_map}
\end{figure}

\textbf{Key finding:} Dense fire hotspots identified as clusters; isolated events classified as noise---useful for distinguishing systematic fire-prone areas from random occurrences.

\section{CLARANS Clustering}

\subsection*{Algorithm Choice}
CLARANS (Clustering Large Applications based on Randomized Search) selected because: (1) uses actual data points as medoids, (2) more robust to outliers than K-Means, (3) more interpretable cluster centers.

\subsection*{Configuration}
\begin{itemize}
    \item Number of clusters: Same as K-Means for comparison
    \item numlocal: 2 (local minima to search)
    \item maxneighbor: Based on dataset size
\end{itemize}

\subsection*{Results}
\begin{figure}[H]
    \centering
    \fbox{\parbox{0.85\textwidth}{\centering\vspace{3cm}\textbf{[PLACEHOLDER: CLARANS Spatial Map with Medoids]}\vspace{3cm}}}
    \caption{CLARANS clusters with medoid locations marked}
    \label{fig:clarans_map}
\end{figure}

\textbf{Advantage:} Medoids are actual fire events, providing interpretable cluster representatives for fire management planning.

\section{Clustering Comparison}

\begin{table}[H]
\centering
\caption{Clustering Methods Comparison}
\label{tab:clustering_comparison}
\begin{tabular}{|l|c|c|c|}
\hline
\textbf{Aspect} & \textbf{K-Means} & \textbf{DBSCAN} & \textbf{CLARANS} \\
\hline
Cluster shape & Spherical & Arbitrary & Spherical \\
Number of clusters & Predefined & Automatic & Predefined \\
Noise handling & No & Yes & No \\
Cluster centers & Centroids & N/A & Medoids \\
Interpretability & Medium & High & High \\
Computational cost & Low & Medium & High \\
\hline
\end{tabular}
\end{table}

\begin{figure}[H]
    \centering
    \fbox{\parbox{0.85\textwidth}{\centering\vspace{4cm}\textbf{[PLACEHOLDER: Side-by-Side Comparison of All Three Methods]}\vspace{4cm}}}
    \caption{Visual comparison of clustering results}
    \label{fig:cluster_comparison}
\end{figure}

\section{Integration with Supervised Learning}
Clustering insights complement classification:
\begin{itemize}
    \item \textbf{Fire risk zones:} Clusters define regions with similar risk profiles
    \item \textbf{Feature validation:} Important features from Random Forest align with cluster-defining features
    \item \textbf{Management application:} Cluster boundaries guide resource allocation for prevention
\end{itemize}

\section{Summary}
Key findings from unsupervised learning:
\begin{enumerate}
    \item \textbf{K-Means:} Identified distinct fire environment zones based on climate and location
    \item \textbf{DBSCAN:} Discovered density-based hotspots; isolated events classified as noise
    \item \textbf{CLARANS:} Provided medoid-based clusters with interpretable centers
\end{enumerate}

\textbf{Practical implications:}
\begin{itemize}
    \item Coastal Mediterranean regions form distinct cluster from interior areas
    \item Climate conditions strongly influence cluster membership
    \item Clustering enables targeted fire prevention by defining risk zones
\end{itemize}

Detailed clustering metrics and parameters are provided in Appendix~\ref{app:clustering}.
